\chapter{svnsync Reference---Subversion Repository
Mirroring}\label{svn.ref.svnsync}

\texttt{svnsync} is the Subversion remote repository mirroring tool. Put
simply, it allows you to replay the revisions of one repository into
another one.

In any mirroring scenario, there are two repositories: the source
repository, and the mirror (or ``sink'') repository. The source
repository is the repository from which \texttt{svnsync} pulls
revisions. The mirror repository is the destination for the revisions
pulled from the source repository. Each of the repositories may be local
or remote---they are only ever addressed by their URLs.

The \texttt{svnsync} process requires only read access to the source
repository; it never attempts to modify it. But obviously,
\texttt{svnsync} requires both read and write access to the mirror
repository.

\texttt{svnsync} is very sensitive to changes made in the mirror
repository that weren\textquotesingle t made as part of a mirroring
operation. To prevent this from happening, it\textquotesingle s best if
the \texttt{svnsync} process is the only process permitted to modify the
mirror repository.

Options in \texttt{svnsync} are global, just as they are in \texttt{svn}
and \texttt{svnadmin}:

\protect\phantomsection\label{svn.ref.svnsync.sw}{}

\begin{description}
\item[\protect\phantomsection\label{svn.ref.svnsync.sw.allow_non_empty}{}\texttt{-\/-allow-non-empty}]
Disables the verification (which \texttt{svnsync\ initialize} performs
by default) that the repository being initialized is empty of history
version.
\item[\protect\phantomsection\label{svn.ref.svnsync.sw.config_dir}{}\texttt{-\/-config-dir}
\textless DIR\textgreater{}]
Instructs Subversion to read configuration information from the
specified directory instead of the default location
(\texttt{.subversion} in the user\textquotesingle s home directory).
\item[\protect\phantomsection\label{svn.ref.svnsync.sw.config_option}{}\texttt{-\/-config-option}
\textless CONFSPEC\textgreater{}]
Sets, for the duration of the command, the value of a runtime
configuration option. \textless CONFSPEC\textgreater{} is a string which
specifies the configuration option namespace, name and value that
you\textquotesingle d like to assign, formatted as
\textless FILE\textgreater:\textless SECTION\textgreater:\textless OPTION\textgreater={[}\textless VALUE\textgreater{]}.
In this syntax, \textless FILE\textgreater{} and
\textless SECTION\textgreater{} are the runtime configuration file
(either \texttt{config} or \texttt{servers}) and the section thereof,
respectively, which contain the option whose value you wish to change.
\textless OPTION\textgreater{} is, of course, the option itself, and
\textless VALUE\textgreater{} the value (if any) you wish to assign to
the option. For example, to temporarily disable the use of the
compression in the HTTP protocol, use
\texttt{-\/-config-option=servers:global:http-compression=no}. You can
use this option multiple times to change multiple option values
simultaneously.
\item[\protect\phantomsection\label{svn.ref.svnsync.sw.disable_locking}{}\texttt{-\/-disable-locking}]
Causes \texttt{svnsync} to bypass its own exclusive access mechanisms
and operate on the assumption that its exclusive access to the mirror
repository is being guaranteed through some other, out-of-band
mechanism.
\item[\protect\phantomsection\label{svn.ref.svnsync.sw.no_auth_cache}{}\texttt{-\/-no-auth-cache}]
Prevents caching of authentication information (e.g., username and
password) in the Subversion runtime configuration directories.
\item[\protect\phantomsection\label{svn.ref.svnsync.sw.non_interactive}{}\texttt{-\/-non-interactive}]
In the case of an authentication failure or insufficient credentials,
prevents prompting for credentials (e.g., username or password). This is
useful if you\textquotesingle re running Subversion inside an automated
script and it\textquotesingle s more appropriate to have Subversion fail
than to prompt for more information.
\item[\protect\phantomsection\label{svn.ref.svnsync.sw.quiet}{}\texttt{-\/-quiet}
(\texttt{-q})]
Requests that the client print only essential information while
performing an operation.
\item[\protect\phantomsection\label{svn.ref.svnsync.sw.revision}{}\texttt{-\/-revision}
(\texttt{-r}) \textless ARG\textgreater{}]
Used by \texttt{svnsync\ copy-revprops} to specify a particular revision
or revision range on which to operate.
\item[\protect\phantomsection\label{svn.ref.svnsync.sw.source_password}{}\texttt{-\/-source-password}
\textless PASSWD\textgreater{}]
Specifies the password for the Subversion server from which you are
syncing. If not provided, or if incorrect, Subversion will prompt you
for this information as needed.
\item[\protect\phantomsection\label{svn.ref.svnsync.sw.source_prop_encoding}{}\texttt{-\/-source-prop-encoding\ ARG}]
Instructs \texttt{svnsync} to assume that translatable Subversion
revision properties found in the source repository are stored using the
character encoding \textless ARG\textgreater{} and to transcode those
into UTF-8 when copying them into the mirror repository.
\item[\protect\phantomsection\label{svn.ref.svnsync.sw.source_username}{}\texttt{-\/-source-username}
\textless NAME\textgreater{}]
Specifies the username for the Subversion server from which you are
syncing. If not provided, or if incorrect, Subversion will prompt you
for this information as needed.
\item[\protect\phantomsection\label{svn.ref.svnsync.sw.steal_lock}{}\texttt{-\/-steal-lock}]
Causes \texttt{svnsync} to steal, as necessary, the lock which it uses
on the mirror repository to ensure exclusive repository access. (This
option should only be used when a lock exists in the mirror repository
and is known to be stale---that is, when you are certain that there are
no other \texttt{svnsync} processes accessing that repository.)
\item[\protect\phantomsection\label{svn.ref.svnsync.sw.sync_password}{}\texttt{-\/-sync-password}
\textless PASSWD\textgreater{}]
Specifies the password for the Subversion server to which you are
syncing. If not provided, or if incorrect, Subversion will prompt you
for this information as needed.
\item[\protect\phantomsection\label{svn.ref.svnsync.sw.sync_username}{}\texttt{-\/-sync-username}
\textless NAME\textgreater{}]
Specifies the username for the Subversion server to which you are
syncing. If not provided, or if incorrect, Subversion will prompt you
for this information as needed.
\item[\protect\phantomsection\label{svn.ref.svnsync.sw.trust_server_cert}{}\texttt{-\/-trust-server-cert}]
Used with \texttt{-\/-non-interactive} to accept any unknown SSL server
certificates without prompting.
\end{description}

\section{svnsync copy-revprops}\label{svn.ref.svnsync.c.copy-revprops}

\emph{Copy all revision properties for a particular revision (or range
of revisions) from the source repository to the mirror repository.}

\textbf{Synopsis}

\texttt{svnsync\ copy-revprops\ DEST\_URL\ {[}SOURCE\_URL{]}}

\texttt{svnsync\ copy-revprops\ DEST\_URL\ REV{[}:REV2{]}}

\subsection{Description}

Because Subversion revision properties can be changed at any time,
it\textquotesingle s possible that the properties for some revision
might be changed after that revision has already been synchronized to
another repository. Because the \texttt{svnsync\ synchronize} command
operates only on the range of revisions that have not yet been
synchronized, it won\textquotesingle t notice a revision property change
outside that range. Left as is, this causes a deviation in the values of
that revision\textquotesingle s properties between the source and mirror
repositories. \texttt{svnsync\ copy-revprops} is the answer to this
problem. Use it to resynchronize the revision properties for a
particular revision or range of revisions.

When \textless SOURCE\_URL\textgreater{} is provided, \texttt{svnsync}
will use it as the repository URL which the destination repository is
mirroring. Generally, \textless SOURCE\_URL\textgreater{} will be
exactly the same source URL as was used with the
\texttt{svnsync\ initialize} command when the mirror was first set up.
You may choose, however, to omit \textless SOURCE\_URL\textgreater, in
which case \texttt{svnsync} will consult the mirror
repository\textquotesingle s records to determine the source URL which
should be used.

We strongly recommend that you specify the source URL on the
command-line, especially when untrusted users have write access to the
revision 0 properties which \texttt{svnsync} uses to coordinate its
efforts.

\subsection{Options}

\begin{verbatim}
--config-dir DIR
--config-option CONFSPEC
--disable-locking
--force-interactive
--no-auth-cache
--non-interactive
--quiet (-q)
--revision (-r) ARG
--skip-unchanged
--source-password PASSWD
--source-prop-encoding ARG
--source-trust-server-cert-failures ARG
--source-username NAME
--steal-lock
--sync-password PASSWD
--sync-trust-server-cert-failures ARG
--sync-username NAME
--trust-server-cert
\end{verbatim}

\subsection{Examples}

Resynchronize the revision properties associated with a single revision
(r6):

\begin{verbatim}
$ svnsync copy-revprops -r 6 file:///var/svn/repos-mirror \
                             http://svn.example.com/repos
Copied properties for revision 6.
$
\end{verbatim}

\section{svnsync help}\label{svn.ref.svnsync.c.help}

\emph{Help!}

\textbf{Synopsis}

\texttt{svnsync\ help}

\subsection{Description}

This subcommand is useful when you\textquotesingle re trapped in a
foreign prison with neither a Net connection nor a copy of this book,
but you do have a local Wi-Fi network running and you\textquotesingle d
like to sync a copy of your repository over to the backup server that
Ira The Knife is running over in cell block D.

\subsection{Options}

None

\section{svnsync info}\label{svn.ref.svnsync.c.info}

\emph{Print information about the synchronization of a destination
repository.}

\textbf{Synopsis}

\texttt{svnsync\ info\ DEST\_URL}

\subsection{Description}

Print the synchronization source URL, source repository UUID and the
last revision merged from the source to the destination repository at
\textless DEST\_URL\textgreater.

\subsection{Options}

\begin{verbatim}
--config-dir DIR
--config-option CONFSPEC
--force-interactive
--no-auth-cache
--non-interactive
--source-password PASSWD
--source-trust-server-cert-failures ARG
--source-username NAME
--sync-password PASSWD
--sync-trust-server-cert-failures ARG
--sync-username NAME
--trust-server-cert
\end{verbatim}

\subsection{Examples}

Print the synchronization information of a mirror repository:

\begin{verbatim}
$ svnsync info file:///var/svn/repos-mirror
Source URL: http://svn.example.com/repos
Source Repository UUID: e7fe1b91-8cd5-0310-98dd-2f12e793c5e8
Last Merged Revision: 47
$
\end{verbatim}

\section{svnsync initialize (init)}\label{svn.ref.svnsync.c.init}

\emph{Initialize a mirror repository for synchronization from the source
repository.}

\textbf{Synopsis}

\texttt{svnsync\ initialize\ MIRROR\_URL\ SOURCE\_URL}

\subsection{Description}

\texttt{svnsync\ initialize} verifies that a repository meets the basic
requirements of a new mirror repository and records the initial
administrative information that associates the mirror repository with
the source repository (specified by \textless SOURCE\_URL\textgreater).
This is the first \texttt{svnsync} operation you run on a would-be
mirror repository.

Ordinarily, \textless SOURCE\_URL\textgreater{} is the URL of the root
directory of the Subversion repository you wish to mirror. Subversion
1.5 and newer allow you to use \texttt{svnsync} for partial repository
mirroring, though --- simply specify the URL of the source repository
subdirectory you wish to mirror as \textless SOURCE\_URL\textgreater.

By default, the aforementioned basic requirements of a mirror are that
it allows revision property modifications and that it contains no
version history. However, as of Subversion 1.7, you may now optionally
disable the verification that the target repository is empty using the
\texttt{-\/-allow-non-empty} option. While the use of this option should
not become habitual (as it bypasses a valuable safeguard mechanism), it
does aid in one very common use-case: initializing a copy of a
repository as a mirror of the original. This is especially handy when
setting up new mirrors of repositories which contain a large amount of
version history. Rather than initialize a brand new repository as a
mirror and then syncronize all of the history into it, administrators
will find it \emph{significantly} faster to first make a copy of the
mature repository (perhaps using \texttt{svnadmin\ hotcopy}) and then
use \texttt{svnsync\ initialize\ -\/-allow-non-empty} to initialize that
copy as a mirror which is now already up-to-date with the original.

\subsection{Options}

\begin{verbatim}
--allow-non-empty
--config-dir DIR
--config-option CONFSPEC
--disable-locking
--force-interactive
--no-auth-cache
--non-interactive
--quiet (-q)
--source-password PASSWD
--source-prop-encoding ARG
--source-trust-server-cert-failures ARG
--source-username NAME
--steal-lock
--sync-password PASSWD
--sync-trust-server-cert-failures ARG
--sync-username NAME
--trust-server-cert
\end{verbatim}

\subsection{Examples}

Fail to initialize a mirror repository due to inability to modify
revision properties:

\begin{verbatim}
$ svnsync initialize file:///var/svn/repos-mirror \
                     http://svn.example.com/repos
svnsync: Repository has not been enabled to accept revision propchanges;
ask the administrator to create a pre-revprop-change hook
$
\end{verbatim}

Initialize a repository as a mirror, having already created a
pre-revprop-change hook that permits all revision property changes:

\begin{verbatim}
$ svnsync initialize file:///var/svn/repos-mirror \
                     http://svn.example.com/repos
Copied properties for revision 0.
$
\end{verbatim}

\section{svnsync synchronize (sync)}\label{svn.ref.svnsync.c.sync}

\emph{Transfer all pending revisions from the source repository to the
mirror repository.}

\textbf{Synopsis}

\texttt{svnsync\ synchronize\ DEST\_URL\ {[}SOURCE\_URL{]}}

\subsection{Description}

The \texttt{svnsync\ synchronize} command does all the heavy lifting of
a repository mirroring operation. After consulting with the mirror
repository to see which revisions have already been copied into it, it
then begins to copy any not-yet-mirrored revisions from the source
repository.

\texttt{svnsync\ synchronize} can be gracefully canceled and restarted.

When \textless SOURCE\_URL\textgreater{} is provided, \texttt{svnsync}
will use it as the repository URL which the destination repository is
mirroring. Generally, \textless SOURCE\_URL\textgreater{} will be
exactly the same source URL as was used with the
\texttt{svnsync\ initialize} command when the mirror was first set up.
You may choose, however, to omit \textless SOURCE\_URL\textgreater, in
which case \texttt{svnsync} will consult the mirror
repository\textquotesingle s records to determine the source URL which
should be used.

We strongly recommend that you specify the source URL on the
command-line, especially when untrusted users have write access to the
revision 0 properties which \texttt{svnsync} uses to coordinate its
efforts.

\subsection{Options}

\begin{verbatim}
--config-dir DIR
--config-option CONFSPEC
--disable-locking
--force-interactive
--no-auth-cache
--non-interactive
--quiet (-q)
--source-password PASSWD
--source-prop-encoding ARG
--source-trust-server-cert-failures ARG
--source-username NAME
--steal-lock
--sync-password PASSWD
--sync-trust-server-cert-failures ARG
--sync-username NAME
--trust-server-cert
\end{verbatim}

\subsection{Examples}

Copy unsynchronized revisions from the source repository to the mirror
repository:

\begin{verbatim}
$ svnsync synchronize file:///var/svn/repos-mirror \
                      http://svn.example.com/repos
Committed revision 1.
Copied properties for revision 1.
Committed revision 2.
Copied properties for revision 2.
Committed revision 3.
Copied properties for revision 3.
…
Committed revision 45.
Copied properties for revision 45.
Committed revision 46.
Copied properties for revision 46.
Committed revision 47.
Copied properties for revision 47.
$
\end{verbatim}

